% Generated draft: requires mathematical and source review.
\documentclass[11pt]{article}
\usepackage[margin=25mm]{geometry}
\usepackage{fontspec}
\setmainfont{Latin Modern Roman}
\usepackage{amsmath,amssymb,mathtools,unicode-math}
\setmathfont{Latin Modern Math}
\usepackage{enumitem,needspace,xurl,hyperref}
\setlength{\parindent}{0pt}
\setlength{\parskip}{5pt}
\setlength{\emergencystretch}{3em}
\newcommand{\N}{\mathbb{N}}
\newcommand{\Z}{\mathbb{Z}}
\newcommand{\Q}{\mathbb{Q}}
\newcommand{\R}{\mathbb{R}}
\newcommand{\C}{\mathbb{C}}
\newcommand{\F}{\mathbb{F}}
\newcommand{\A}{\mathbb{A}}
\newcommand{\PP}{\mathbb{P}}
\newcommand{\E}{\mathbb{E}}
\newcommand{\eps}{\varepsilon}
\newcommand{\dd}{\,\mathrm{d}}
\begin{document}
\section*{UnsolvedMath ID: 20001986}
\textbf{Source problem code:} catalogue entry
\subsection{Definitions and mathematical statement}
Einsiedler and Lind ask for further recognition conditions ensuring that a totally concave polyhedral complex is the nonarchimedean amoeba of a prime Laurent ideal. Let \(k\) be algebraically closed with a nonarchimedean valuation \(\nu:k^*\to\Gamma\subset\R\), and let \(\mathfrak p\subset k[x_1^{\pm1},\ldots,x_d^{\pm1}]\) be prime. In the nontrivially valued setting, write
\[
 \operatorname{Trop}(\mathfrak p)=
 \overline{\{(\nu(x_1),\ldots,\nu(x_d)):
                  x\in V(\mathfrak p)(k)\}}\subset\R^d.
\]
The background identifies this with the appropriate Bieri--Groves set, of pure dimension equal to the quotient algebra's Krull dimension, rationally defined over the value group. Here lattice-rational directions and value-group-defined offsets are distinct requirements.

The source calls the complex homogeneous. If this means a fan invariant under positive scaling about the origin, that is an additional restriction: tropicalizations with nontrivially valued coefficients need not be homogeneous fans. The aid's finite triangle example below is not homogeneous in that sense. Thus it gives an obstruction for the more general polyhedral recognition problem, not directly a counterexample satisfying an unrecovered fan hypothesis.

For a finite pure lattice-rational \(m\)-dimensional complex \(\Sigma\), choose the ambient lattice \(N=\Z^d\). At a codimension-one face \(\tau\), let \(N_\tau=N\cap\operatorname{span}(\tau-\tau)\). For every incident maximal cell \(\sigma\), let \(u_{\sigma/\tau}\) be the primitive outward direction in the quotient lattice. Positive integral tropical multiplicities must satisfy
\[
 \sum_{\sigma\supset\tau}w_\sigma u_{\sigma/\tau}=0
 \quad\text{in }(N/N_\tau)\otimes\R,
 \qquad w_\sigma\in\Z_{>0}.
\]
Choosing quotient bases assembles these equations into an integer matrix \(B_\Sigma\). The aid's necessary condition is strict positive-kernel feasibility:
\[
 \exists w\in\R_{>0}^{\Sigma^{(m)}}:B_\Sigma w=0.
\]
For an integer matrix this is equivalent to positive rational and positive integral feasibility. Stiemke's alternative gives an infeasibility certificate
\[
 \exists y:\quad B_\Sigma^Ty\geq0,
 \qquad B_\Sigma^Ty\ne0.
\]
Such a certificate is incompatible with \(w>0\), by pairing with \(w\). Balancing is necessary, not sufficient for prime tropical realizability.

The explicit one-dimensional example has vertices
\[
 a=(0,0),\quad b=(1,0),\quad c=(0,1),
\]
all three triangle edges, and rays based at \(a,b,c\) in respective primitive directions
\[
 (-1,-1),\qquad(2,-1),\qquad(-1,3).
\]
It is pure and connected through codimension one. Locally at each vertex its three outgoing directions have a positive dependence: at \(a\) the weights on \(ab,ac\) and its ray are equal; at \(b\) those on \(ab,bc\) and its ray are equal; at \(c\), the \(ac\) weight is twice the equal \(bc\) and ray weights. These local stars satisfy the one-dimensional total-concavity property used by the aid: a locally supporting half-space cannot contain all directions unless its boundary contains the whole local star. Edge-interior stars satisfy it as well.

Globally, if \(A,B,C\) are the weights of \(ab,ac,bc\), the vertex equations force
\[
 A=B,\qquad A=C,\qquad B=2C.
\]
Hence every balanced weighting is zero, including the ray weights. The local positive ratios accumulate multiplicative holonomy two around the triangle and cannot be made into global multiplicities. This excludes the support from being a prime tropicalization, indeed from supporting any nonzero balanced tropical curve. The coordinates are integral; over a value group not containing one the construction can be rescaled by a positive group element without changing its obstruction.

The aid uses this to show that local concavity and codimension-one connectedness do not replace global positive balancing. The source's precise higher-dimensional total-concavity definition, homogeneous restriction and complete necessary-and-sufficient recognition target remain unrecovered; no general realization theorem follows.
\subsection{Short English statement}
A prime tropicalization needs globally compatible positive integral balancing weights, not just local concavity. The aid's triangle-with-rays has positive balances at every vertex but incompatible ratios around its cycle. Its nonconical example does not settle a version restricted to homogeneous fans.
\subsection{Sources}
\begin{itemize}
\item Permanent UnsolvedMath ID \(20001986\), exact public code catalogue entry; local registry \texttt{aim-geometry-notes.json}, \texttt{notes[323]}, canonical tag section.
\item AIM workshop ``Amoebas and tropical geometry,'' item A.4 ``Recognition problems,'' contributed by Manfred Einsiedler and Doug Lind: \url{https://aimath.org/WWN/amoebas/amoebas.pdf}. No PDF was accessed.
\item Original background names the Bieri--Groves set and total concavity without a full definition or bibliography. The aid invokes integer balancing and Stiemke's alternative, also without full references. No original literature entry or year is supplied.
\item Research attempt 1, ``Global balancing as a recognition obstruction,'' labelled partial\_result, open at the run, proved and low novelty confidence; registry partially\_solved. The exact-statement verification note explicitly describes a damaged extraction and itself ends mid-sentence.
\end{itemize}
\end{document}
